%!TEX root = ../main/main.tex

En esta pregunta usamos la siguiente notación que será utilizada en el
curso. Dados dos strings $u, v$, el string $u \| v$ se define como la
concatenación de $u$ con $v$. Además, dados dos strings binarios $u,
v$ del mismo largo, $u \oplus v$ es el string binario que resulta de
tomar el XOR de $u$ con $v$ componente a componente (lo cual
corresponde a tomar suma en módulo 2). Por ejemplo, si $u = 00100$ y
$v = 11110$, entonces $u \| v = 0010011110$ y $u \oplus v = 11010$.


Sea $n$ un número natural, y suponga que $\{f_k\}_{k \in \{0,1\}^n}$
es una familia de funciones tales que $f_k
: \{0,1\}^n \to \{0,1\}^n$. Note que estas funciones no son
necesariamente biyectivas. A partir de esta familia definimos un
esquema criptográfico $(\textit{Gen}, \textit{Enc}, \textit{Dec})$
sobre los espacios $\mathcal{M} = \mathcal{C} = \mathcal{K}
= \{0,1\}^{2n}$. En particular, la función de encriptación
$\textit{Enc}$ es definida de la siguiente forma considerando un
mensaje $m \in \{0,1\}^{2n}$ tal que $m = m_1 \| m_2$ con $m_1,
m_2 \in \{0,1\}^n$, y una clave $k \in \{0,1\}^{2n}$ tal que $k =
k_1 \| k_2$ con $k_1, k_2 \in \{0,1\}^n$:
\begin{align*}
\textit{Enc}_k(m) \ = \ (m_1 \oplus f_{k_1}(m_2)) \| (m_2 \oplus f_{k_2}(m_1 \oplus f_{k_1}(m_2))). 
\end{align*}
En este problema usted debe responder las siguientes preguntas.
\begin{enumerate}
\item {[2 puntos]} Indique como se define la función $\textit{Dec}$. Es decir, dado una clave $k \in \{0,1\}^{2n}$ y un mensaje cifrado $c \in \{0,1\}^{2n}$, indique cómo se calcula el texto plano $m$ tal que $\textit{Enc}_k(m) = c$.

\item {[4 puntos]} Demuestre que el esquema $(\textit{Gen}, \textit{Enc}, \textit{Dec})$ no es una pseudo-random permutation (PRP) con 2 rondas.  
\end{enumerate}
Note que (b) muestra que $(\textit{Gen}, \textit{Enc}, \textit{Dec})$
no es una PRP independientemente de las propiedades que cumpla la
familia de funciones $\{f_k\}$.


\paragraph{Corrección.} La asignación de puntajes para cada pregunta de este problema es la siguiente.
\begin{itemize}
\item[(a)] En esta pregunta se obtiene 1 punto si se indica cómo se define la función $\textit{Dec}$, pero la solución está incompleta o tiene algunos errores menores. Y se obtiene 2 puntos si se indica cómo se define la función $\textit{Dec}$, y la descripción entregada es correcta. 

\item[(b)] Esta pregunta se corrige considerando que se debe definir la
estrategia del adversario que le permite ganar el juego de $2$ rondas
que define una PRP con una probabilidad mayor o igual a $\frac{1}{2} +
f(n)$, donde $f(n)$ es una función no despreciable. La asignación de
puntaje en esta pregunta es la siguiente.
\begin{itemize}
    \item{[1 puntos]} Sólo se entrega la definición de la estrategia
    del adversario. Esta estrategia tiene elementos que efectivamente
    podrían permitir al adversario ganar con una probabilidad
    significativamente mayor a $\frac{1}{2}$, pero tiene algunos
    errores.

    \item{[2 puntos]} Se entrega una definición de la estrategia del
    adversario que está correcta en el sentido de que le permite al
    adversario ganar el juego con probabilidad mayor o igual a
    $\frac{1}{2} + f(n)$, donde $f(n)$ es una función no despreciable.

    \item{[3 puntos]} Se entrega una definición de la estrategia del
    adversario que está correcta en el sentido de que le permite al
    adversario ganar el juego con probabilidad mayor o igual a
    $\frac{1}{2} + f(n)$, donde $f(n)$ es una función no
    despreciable. Además se explica por qué la probabilidad de que el
    adversario gane el juego es mayor o igual a $\frac{1}{2}
    + f(n)$, pero no se entrega una demostración completa de
    esta propiedad.

    \item{[4 puntos]} Se entrega una definición de la estrategia del
    adversario que está correcta en el sentido de que le permite al
    adversario ganar el juego con probabilidad mayor o igual a
    $\frac{1}{2} + f(n)$, donde $f(n)$ es una función no
    despreciable. Además se entrega una demostración
    formal de que la probabilidad de que el adversario gane el juego
    es mayor o igual a $\frac{1}{2} + f(n)$.
\end{itemize}
\end{itemize}

